Ein Bezugssystem ist genau dann ein Inertialsystem, wenn es sich gegenüber einem Inertialsystem geradlinig gleichförmig bewegt (oder ruht), also nicht beschleunigt ist. Nur dann gilt in ihm das Trägheitsgesetz.
\begin{abcliste}
\abc Das Auto ändert in der Kurve die Richtung seiner Geschwindigkeit, es ist beschleunigt: $\result{\text{kein Inertialsystem}}$. (Gegenstände im Auto rutschen ohne erkennbare Ursache nach aussen.)
\abc Das Pendel ändert dauernd Betrag und Richtung seiner Geschwindigkeit: $\result{\text{kein Inertialsystem}}$.
\abc Der Zug bewegt sich gegenüber der Erdoberfläche geradlinig gleichförmig: $\result{\text{Inertialsystem}}$.
\end{abcliste}
